%----------------------------------------------------------------------------------------
%	PAQUETES Y OTRAS CONFIGURACIONES
%----------------------------------------------------------------------------------------

%----------------------------------------------------------------------------------------
%	PAQUETES Y OTRAS CONFIGURACIONES
%----------------------------------------------------------------------------------------

\documentclass[paper=letter, fontsize=11pt]{scrartcl} % Tamaño de papel y letra para el documento

\usepackage[utf8]{inputenc} % Los caracteres acentuados se pueden escribir normalmente en el código
\usepackage[T1]{fontenc} % Configuración de fuente de salida
\usepackage{fourier} % Se usa una fuente diferente al default
\usepackage[spanish,es-noquoting]{babel} % Se configura como documento en español
\usepackage{amsmath,amsfonts,amsthm} % Paquetes para escribir formulas matemáticas
\usepackage{graphicx} % Paquetes para incluir imágenes

\usepackage{circuitikz}
\usepackage{tikz}
\usetikzlibrary{arrows}

\usepackage{sectsty} % Paquete para configuración de secciones
\allsectionsfont{\centering \normalfont \scshape} % Los títulos de las secciones son centrados, con la misma fuente y pequeñas mayúsculas

\usepackage{todonotes}
\usepackage{microtype}

\usepackage{fancyhdr} % Paquete para personalizar pies y cabeceras de página
\pagestyle{fancyplain} % Todas las páginas con las mismas cabeceras y pies de página
\fancyhead{} % Sin cabecera
\fancyfoot[L]{} % Vacío en la izquierda del pie de página
\fancyfoot[C]{} % Vacío en el centro del pie de página
\fancyfoot[R]{\thepage} % Número de página en el pie de pagina
\renewcommand{\headrulewidth}{0pt} % Sin lineas en la cabecera
\renewcommand{\footrulewidth}{0pt} % Sin lineas en el pie de página
\setlength{\headheight}{13.6pt} % Altura de cabecera

\numberwithin{equation}{section} % Numera ecuaciones en cada sección
\numberwithin{figure}{section} % Numera figuras en cada sección
\numberwithin{table}{section} % Numera tablas en cada sección

\setlength\parindent{0pt} % Quita la indentación de los párrafos

\newcommand{\horrule}[1]{\rule{\linewidth}{#1}} % Comando personalizado para hacer linea horizontal


%----------------------------------------------------------------------------------------
%	TITULO
%----------------------------------------------------------------------------------------

\title{	
	\normalfont \normalsize
	\begin{figure}[h]
		\begin{center}
			\includegraphics[width=0.3\textwidth]{../../images/UNITEC.png} % La imagen esta incluida en el mismo directorio del código
		\end{center}
	\end{figure}
	\textsc{Instrumentación y Control} \\ [25pt]
	\horrule{0.5pt} \\[0.4cm] % Linea horizontal delgada
	\huge Práctica 3 - Medición de variables físicas II \\ % Titulo de la práctica
	\horrule{2pt} \\[0.5cm] % Linea horizontal mas gruesa
}

\author{Roberto Cadena Vega} % Nombre del profesor

\date{\normalsize 20 de febrero de 2014} % Fecha de la práctica

%----------------------------------------------------------------------------------------
%	EMPIEZA EL DOCUMENTO
%----------------------------------------------------------------------------------------

\begin{document}

\maketitle % Imprime el título

%----------------------------------------------------------------------------------------
%	OBJETIVOS
%----------------------------------------------------------------------------------------

\section{Objetivos}

	Implementar un sistema eléctrico simple, que sea capaz de medir la intensidad luminosa.

%----------------------------------------------------------------------------------------
%	CONOCIMIENTOS PREVIOS
%----------------------------------------------------------------------------------------

\section{Conocimientos Previos}

%----------------------------------------------------------------------------------------

	\subsection{Transductores}

		El día de hoy usaremos un sensor de fotones (es decir, de luz). Estrictamente hablando, este es un transductor solamente, pero podemos utilizarlo para medir una variable física simplemente.

		La fotorresistencia que utilizaremos es básicamente una resistencia, que varía su valor dependiendo de la intensidad luminosa. Por lo general las venden en valores como $1 k \Omega$ o $10 k \Omega$, y esto quiere decir que variarán desde $0 \Omega$, hasta el valor nominal.

		El problema es que la resistencia no es una señal que podamos medir fácilmente, de hecho cuando la medimos con el multímetro, este tiene que aplicar un voltaje de prueba para saber, a través de un cálculo, la resistencia y justo esto es lo que haremos.

%----------------------------------------------------------------------------------------

	\subsection{Acondicionamiento}

		En esta ocasión no queremos una señal analógica como salida de nuestro sistema, sino una señal que nos diga simplemente si hay luz o no, por lo que no podemos simplemente amplificar la señal que mediremos en nuestra resistencia.

		\begin{center}
			\tikzstyle{input} = [coordinate]
			\tikzstyle{output} = [coordinate]
			\tikzstyle{block} = [draw, rectangle, minimum height=3em, minimum width=6em]
			\tikzstyle{init} = [pin edge={to-, thin, black}]

			\begin{tikzpicture}[auto, node distance=4cm, >=latex']
				\node [input, name=entrada] {};
				\node [block, right of=entrada] (sensor) {Señal};
				\node [block, right of=sensor] (acondicionador) {Acondicionamiento};
				\node [output, right of=acondicionador] (salida) {};

				\draw [->] (entrada) -- node[name=x] {$x$} (sensor);
				\draw [->] (sensor) -- node[name=x] {$u$} (acondicionador);
				\draw [->] (acondicionador) -- node[name=y] {$y$} (salida);
			\end{tikzpicture}
		\end{center}

		¿Más aún, qué tan oscuro significa que no hay luz? ¿Si alguien saca su celular e ilumina ligeramente el cuarto debemos considerar que hay luz? Tenemos que considerar una referencia.

%----------------------------------------------------------------------------------------

		\subsubsection{Comparadores}

			Los OPAMP's que usamos en la práctica anterior tambien pueden hacer una comparación entre dos señales (de hecho, así es como logran amplificar una señal, pero por ahora no entraremos en detalle con eso), por lo que ahora revisaremos cómo se puede configurar un OPAMP para que funcione como comparador.

			\begin{center}
				\begin{circuitikz} \draw
					(0,0) node[op amp](opamp){}
					(opamp.+) to [short, -*] (-1.2, -0.5) node[below] {$V_i$}
					(opamp.-) to [short, -o] (-1.2, 0.5) node[left] {$V_{ref}$}
					(opamp.out) to [short, -o] (1.5, 0) node[right] {$V_o$}
					(opamp.down) node[ground] {}
					(opamp.up) to [short, -o] (-0.08, 1) node[above] {$V_+$}
					(-1.2, -0.5) to [short, *-] (-3, -0.5)
					(-3, -0.5) to [vR, l=$R_F$] (-3, 2.5) node[above] {$5 V$}
					(-3, -0.5) to [R, l=$R_1$] (-3, -3.5) node[ground] {} 
				;\end{circuitikz}	
			\end{center}
			
			y la ecuación que rige su comportamiento es:

				\begin{equation}
					V_o = \left\{
					\begin{array}{r l}	
						0 V & V_i < V_{ref}  \\
						V_+ & V_i > V_{ref}  
					\end{array} \right.
				\end{equation}

			eso quiere decir que el diagrama queda como sigue:

			\begin{center}
				\tikzstyle{input} = [coordinate]
				\tikzstyle{output} = [coordinate]
				\tikzstyle{block} = [draw, rectangle, minimum height=3em, minimum width=6em]
				\tikzstyle{sum} = [draw, circle, node distance=3cm]
				\tikzstyle{init} = [pin edge={to-, thin, black}]
				\tikzstyle{pinstyle} = [pin edge={to-,thin,black}]

				\begin{tikzpicture}[auto, node distance=3cm, >=latex']
					\node [input, name=entrada] {};
					\node [block, right of=entrada] (sensor) {Señal};
					\node [sum, right of=sensor, pin={[pinstyle]above:$- V_{ref}$}] (suma) {};
					\node [block, right of=suma] (amplificador) {$\infty$};
					\node [output, right of=amplificador] (salida) {};

					\draw [->] (entrada) -- node {$x$} (sensor);
					\draw [->] (sensor) -- node {$u$} node[pos=0.99] {$+$} (suma);
					\draw [->] (suma) -- node {$v$} (amplificador);
					\draw [->] (amplificador) -- node {$y$} (salida);
				\end{tikzpicture}
			\end{center}

			Por ahora no te preocupes por el $\infty$, tan solo hay que notar que nuestra salida tendrá un valor de $V_+$ o $0 V$, dependiendo de la entrada.

			Recuerda que dependiendo del OPAMP que uses, las terminales de tu circuito integrado serán diferentes, por lo que es absolutamente necesario que tengas el datasheet de tu OPAMP a la mano.

%----------------------------------------------------------------------------------------
%	EQUIPO
%----------------------------------------------------------------------------------------

\section{Equipo}

	El siguiente equipo será proporcionado por el laboratorio, siempre y cuando lleguen en los primeros 15 minutos de la práctica, y hagan el vale conteniendo el siguiente equipo (exceptuando las pinzas).

	\begin{itemize}
		\item 1 Fuente de Alimentación
		\item 1 Multímetro
		\item 1 Cable de alimentación
		\item 2 Cables banana - caimán
		\item Pinzas
	\end{itemize}

%----------------------------------------------------------------------------------------
%	MATERIALES
%----------------------------------------------------------------------------------------

\section{Materiales}

	\begin{itemize}
		\item Protoboard
		\item Fotorresistencia de $1k \Omega$ o $10 k \Omega$ o cualquier otra que consigas
		\item LM741 o LM358
		\item Resistencias (considera una resistencia del mismo valor que la fotorresistencia)
		\begin{itemize}
			\item $180 \Omega$
			\item $220 \Omega$
			\item $330 \Omega$
			\item $1 k\Omega$
			\item $3.3 k\Omega$
			\item $10 k\Omega$
		\end{itemize}
		\item Cables
	\end{itemize}

%----------------------------------------------------------------------------------------
%	DESARROLLO
%----------------------------------------------------------------------------------------

\section{Desarrollo}
	
	\begin{enumerate}
		\item Diseña un circuito para medir el voltaje que pasa por la fotorresistencia en serie con otra resistencia (como en la práctica 1, con el LED).
		\item Diseña un circuito para dar un voltaje de referencia al comparador.
		\item Implementa el circuito del OPAMP configurado como comparador.
		\item Realiza las mediciones requeridas en la hoja de anotaciones.
	\end{enumerate}


%----------------------------------------------------------------------------------------
%	CONCLUSIONES
%----------------------------------------------------------------------------------------

\section{Conclusiones}

	El alumno deberá describir sus conclusiones al final de su reporte de práctica.
    
%----------------------------------------------------------------------------------------
%	HOJA DE ANOTACIONES
%----------------------------------------------------------------------------------------

\clearpage
\section{Hoja de Anotaciones}
	
	\begin{enumerate}
		\item ¿Cuál es el valor del voltaje medido en la fotorresistencia cuando hay luz sobre ella? \newline
		\item ¿Cuál es el valor del voltaje medido en la fotorresistencia cuando no hay luz sobre ella? \newline
		\item ¿Cuál es el valor del voltaje medido en la fotorresistencia cuando está completamente cubierta? \newline
		\item ¿Qué valor elegirás como $V_{ref}$? \newline
		\item ¿Qué valores de resistencia elegirás para que el divisor de voltaje te dé $V_{ref}$? \newline
		\item Dibuja el diagrama de todo tu sistema. \newline \newline \newline \newline \newline \newline \newline \newline \newline \newline \newline \newline \newline \newline \newline
	\end{enumerate}

	Integrantes del equipo: \\[0.2cm]
	\horrule{0.5pt} \\[0.2cm] % Linea horizontal delgada
	\horrule{0.5pt} \\[0.2cm] % Linea horizontal delgada
	\horrule{0.5pt} \\[0.2cm] % Linea horizontal delgada
	\horrule{0.5pt} % Linea horizontal delgada	

	Revisó: \\[0.2cm]
	\horrule{0.5pt} \\% Linea horizontal delgada
    
%----------------------------------------------------------------------------------------
%	FIN DEL DOCUMENTO
%----------------------------------------------------------------------------------------

\end{document}